\documentclass[
  11pt,
  a4paper,
  parskip=half         % KOMA-native Absatzabstand (kein parskip-Paket nötig)
]{scrartcl}

% Encoding & Language
\usepackage[utf8]{inputenc}
\usepackage[T1]{fontenc}
\usepackage[ngerman]{babel}

% Fonts: Libertinus-Familie (Serif/Sans/Mono aufeinander abgestimmt)
\usepackage{libertinus}
\usepackage[scaled=0.85]{inconsolata}

% Layout
\usepackage[margin=2.5cm, top=2.8cm, bottom=2.8cm, headheight=22pt]{geometry}
\usepackage{microtype}

% Color
\usepackage[dvipsnames]{xcolor}
\definecolor{accent}{HTML}{1C3A5E}
\definecolor{muted}{HTML}{5A5A5A}

% Section headings: KOMA-native (kein titlesec-Konflikt)
\setkomafont{section}{\sffamily\large\bfseries\color{accent}}
\RedeclareSectionCommand[beforeskip=1.6em, afterskip=0.5em]{section}

% Linie unter section-Überschriften via sectionlinesformat
\makeatletter
\renewcommand{\sectionlinesformat}[4]{%
  \Ifstr{#1}{section}{%
    \@hangfrom{\hskip #2#3}{#4}%
    \par\nobreak\vspace{3pt}%
    {\color{accent!35}\rule{\textwidth}{0.55pt}}%
    \vspace{1pt}%
  }{%
    \@hangfrom{\hskip #2#3}{#4}%
  }%
}
\makeatother

% Header / Footer: KOMA-native (kein fancyhdr-Konflikt)
\usepackage{scrlayer-scrpage}
\clearpairofpagestyles
\ihead{\small\sffamily\textcolor{muted}{Wissenschaftliches Analyseprotokoll}}
\ohead{\small\sffamily\textcolor{muted}{\today}}
\cfoot{\small\sffamily\textcolor{muted}{\pagemark}}
\KOMAoptions{headsepline=0.4pt}
\addtokomafont{headsepline}{\color{accent!30}}

% Lists
\usepackage{enumitem}
\setlist[itemize]{itemsep=0.3em, parsep=0pt, topsep=0.4em, leftmargin=1.4em}
\setlist[itemize,2]{itemsep=0.2em, topsep=0.2em, leftmargin=1.2em}

\setlength{\emergencystretch}{1em}

% ---------------------------------------------------------------
\begin{document}

% Manueller Titel
\begin{center}
  {\sffamily\LARGE\bfseries\color{accent}
    Wissenschaftliches Analyseprotokoll}\\[0.5em]
  {\sffamily\large\color{muted}
    Evaluation des Einleitungs-Entwurfs (Erste Version)}\\[1.2em]
  {\sffamily\small\color{muted}
    Maximilian Wagner\quad\textperiodcentered\quad\today}
\end{center}

\vspace{0.6em}
\par\begingroup\color{accent!35}\hrule height 0.55pt\endgroup
\vspace{0.8em}

% Metadaten-Block
\begin{description}[
  labelwidth=5cm,
  leftmargin=!,
  labelsep=0.6em,
  font=\normalfont\bfseries,
  itemsep=0.25em
]
  \item[Untersuchungsgegenstand] Erste KI-generierte Rohfassung der Einleitung
  \item[Referenzdokument] Exposé (,,Untersuchung mikroarchitektonischer
    Ressourcen-Interfe\-renzen\ldots'')
  \item[Prüfkriterien] Objektivität, Belegprüfung, logische Struktur
\end{description}

\par\begingroup\color{accent!35}\hrule height 0.55pt\endgroup

% ---------------------------------------------------------------
\section{Objektivität und Tonalität}

\begin{itemize}
  \item \textbf{Stärken}
  \begin{itemize}
    \item \textbf{Fachsprache:} Die Tonalität ist durchgehend sachlich und wissenschaftlich.
          Fachbegriffe wie \textit{Last Level Cache}, \textit{Hypervisor}, \textit{CPU-Pinning}
          und \textit{Hardware-Mitigierung} werden im korrekten Kontext und ohne unnötige
          Metaphern verwendet.
    \item \textbf{Faktentrennung:} Die sprachliche Trennung zwischen etablierter Architektur
          (Multi-Tenancy) und der zu prüfenden Hypothese (Isolationsbruch) wird korrekt vollzogen.
  \end{itemize}

  \item \textbf{Schwächen}
  \begin{itemize}
    \item \textbf{Titel-Diskrepanz:} Der vom LLM generierte Titel (,,\textit{Trade-offs der
          Virtualisierung auf x86-Systemen: Eine quantitative Analyse von Performance-Overheads
          und Side-Channel-Risiken}'') weicht inhaltlich stark vom Exposé ab. Er wirkt reißerisch
          und verschiebt den Fokus inkorrekt in Richtung Side-Channels, obwohl die
          Leistungsdegradation (DoS) im Fokus steht. Der präzisere Originaltitel
          (,,\textit{Untersuchung mikroarchitektonischer Ressourcen-Interferenzen (Noisy-Neighbor)
          in virtualisierten Umgebungen}'') muss wiederhergestellt werden.
    \item \textbf{Sprachlicher Stil:} Die Formulierung wirkt stellenweise unnatürlich, gezwungen
          intellektuell und künstlich (KI-typische Phrasierung), was den Lesefluss stört. Die
          Sprache muss natürlicher und neutraler gestaltet werden.
  \end{itemize}
\end{itemize}

% ---------------------------------------------------------------
\section{Belegprüfung und Literaturfokus}

\begin{itemize}
  \item \textbf{Stärken}
  \begin{itemize}
    \item \textbf{Zitations-Format:} Die syntaktische Vorbereitung der Platzhalter
          (z.\,B.\ \texttt{\textbackslash cite\{koh2007analysis\}}) ist für die spätere
          \LaTeX-Transformation und Kompilierung formal korrekt umgesetzt.
  \end{itemize}

  \item \textbf{Schwächen}
  \begin{itemize}
    \item \textbf{Überfokussierung auf Vorab-Literatur:} Der Text baut seine Argumentation
          namentlich auf vorläufigen Quellen auf (,,Wie bereits Koh et~al.\ darlegten\ldots'',
          ,,Die Systematisierung durch Ge et~al.\ldots''). Da die finale Literaturrecherche noch
          aussteht, stellt dies einen strukturellen Fehler dar. In der Einleitung muss das
          Problem allgemein formuliert werden, wobei Quellen lediglich als Beleg in Klammern
          dienen. Die namentliche Diskussion der Literatur gehört ausschließlich in das Kapitel
          ,,Stand der Forschung''.
  \end{itemize}
\end{itemize}

% ---------------------------------------------------------------
\section{Logische Struktur und inhaltliche Tiefe}

\begin{itemize}
  \item \textbf{Stärken}
  \begin{itemize}
    \item \textbf{Der Rote Faden:} Die Einleitung folgt einem kohärenten Trichter-Prinzip
          (deduktiver Aufbau). Sie beginnt breit bei allgemeinen Cloud-Paradigmen (Multi-Tenancy),
          verengt sich auf die mikroarchitektonische Ebene der LLC-Ressourcenteilung und mündet
          präzise in der Forschungsfrage.
    \item \textbf{Methodische Transparenz:} Der Proof of Concept (inklusive der Werkzeuge
          \texttt{stress-ng}, \texttt{sysbench}, \texttt{fio}) sowie die essenzielle
          Fallback-Strategie sind unmissverständlich und detailliert dargelegt und entsprechen
          den Vorgaben des Exposés.
  \end{itemize}

  \item \textbf{Schwächen}
  \begin{itemize}
    \item \textbf{Technologische Überfokussierung:} Die Einleitung verliert sich in einer zu
          hohen technologischen Granularität. Die Nennung spezifischer Angriffsvektoren
          (wie \textsc{Prime+Probe}) ordnet sich nicht der konzeptionellen Abstraktionsebene
          einer Einleitung unter. Die Argumentation muss sich stärker an die übergeordneten
          Konzepte der logischen Isolation (gemäß Exposé) halten.
    \item \textbf{Formale Artefakte:} Das Dokument enthält strukturelle Überreste des verwendeten
          Prompts (z.\,B.\ Aufzählungszeichen wie ,,\texttt{a) Thematische Einführung:}'',
          ,,\texttt{b) Problemstellung \& Relevanz:}''). Eine wissenschaftliche Einleitung im
          IEEE-Format erfordert einen zusammenhängenden Fließtext, der diese Punkte inhaltlich,
          aber nicht als explizite Zwischenüberschriften abarbeitet.
    \item \textbf{Fehlende Abgrenzung zur Basisaufgabe:} Die Fallback-Strategie wird erwähnt,
          der methodische Übergang ist jedoch logisch unzureichend verknüpft. Es muss deutlicher
          formuliert werden, dass der Vergleich der Virtualisierungslösungen (Basisaufgabe) das
          Fundament bildet, auf dem der Noisy-Neighbor-Angriff anschließend evaluiert wird.
  \end{itemize}
\end{itemize}

\end{document}
